\section{Denominator-free and centered three-root transport}
\label{sec:pell-three-root}

Two new eleven-positive-variable endpoints improve the inspected
degree-$325$, $113$-operation construction of source~70:
\[
 \boxed{\deg P_{11}=277,\quad 111\text{ operations}},\qquad
 \boxed{\deg \widetilde P_{11}=257,\quad 114\text{ operations}}.
\]
The first deletes a denominator guard using the adjacent-coefficient
Pell descent of source~69.  The second additionally centers the auxiliary
root and proves transport positivity by a new signed bootstrap.  They are
different degree/cost choices.  Every multiplication by a numeral and every
square is charged; fixed numeral leaves and register reuse are free.
Neither construction changes the known nine-positive-variable bound or
claims global optimality.  All native compiler and literal
machine/refutation calibrations remain named inherited dependencies.

\subsection{A denominator guard that can be deleted}

The following elementary fact replaces the prime-support argument in
source~70.

\paragraph{Adjacent-coefficient descent.}
If $Q>1$ and $V$ are integers, $y$ is rational, and
\[
 QV^2-(Q-1)y^2=1,
\]
then $y$ is an integer.  Indeed, replace both coordinates by their
absolute values and write $y=a/m$ in lowest terms.  Then $m^2\mid Q-1$.
For $V>1$ apply
\[
 V'=(2Q-1)V-2(Q-1)y,\qquad
 y'=(2Q-1)y-2QV.
\]
Here $V'$ is integral, the same norm is preserved, and $y>V$ and
$QV>(Q-1)y$ give $-V<V'<V$.  The norm excludes $V'=0$.
Since $2Q-1\equiv1\pmod m$, the reduced denominator of $y'$ is still
$m$.  Strong induction reaches $V=1$, where $y=1$, and proves $m=1$.

Thus, after a sum-of-three-roots argument has made the auxiliary root
rational, no denominator constraint involving the first transport modulus
is needed.  In particular source~70's $I=iq(q-1)c^2$ can be replaced by
$I=ic^2$, deleting two multiplication gates.  The unconditional separation
proof used only $I\ge c^2$ and remains valid.

\subsection{Complete polynomials}

The eleven unknowns are
\[
 j,F,\alpha,i,T,X,Y,n,\eta,\zeta,\gamma\in\mathbb Z_{>0}.
\]
The four fixed numeral ports are the complete compact44 tuple
$\beta,K_0,M_C,M_F$, with $\beta=2^{590069}-1$ and $M_F$ already
shifted.  They are not free parameters; their authenticated payloads and
whole-tuple binding accompany the circuits.  Every other symbol below is
a polynomial register, not another quantified variable.  Define
\begin{gather*}
 J=3j+1,\qquad q=\beta J+1,\qquad C=q-F-\alpha,\\
 E=XY,\qquad k=\eta+\zeta,\qquad c=kY+\eta,\\
 a=E+Y,\qquad H=4a+3,\qquad \Delta=a^2+H=(a+2)^2-1,\\
 \ell=EkY,\qquad I=ic^2,\qquad M=q^3(q-1),\\
 R=((q-1)(q-F)+\alpha)(q^2-1)+(M_C+qM_F)J,\\
 \widehat R=q^3(q+F+\alpha),\qquad B_t=(qK_0+X)C-qF,\\
 D_\tau=\ell(\ell+k)+1,\qquad D_f=1+\Delta I^2,\\
 Q=\Delta^2I^2,\qquad \delta=Q-1,\qquad V=D_f(cT-R)-c,\\
 r_0=k-R-1,\qquad S=4qY,\qquad U=qS,\\
 D_*=U(X+ac)+H(U\gamma+SX-r_0),\\
 A_1=D_*^2-U^2(1+\Delta c^2).
\end{gather*}

For the degree-$277$ endpoint choose
\[
 L=3D_f(cT+\widehat R),\quad h_*=L,\quad
 W=1+(q-1)h_*,\quad
 Z=M(\ell+k+n+L)+(q-1)Y-q^3WB_t.
\]
For the degree-$257$ endpoint instead choose
\begin{align*}
 B_0&=cT+\widehat R,&L&=D_f+B_0,\\
 h_*&=4qD_f+L,&W&=1+(q-1)h_*,\\
 Z&=M(\ell+k+n+L+V)+(q-1)Y-q^3WB_t.
\end{align*}
The shared division-free final norm is
\begin{align*}
 A&=M^2D_\tau,&B_f&=M^2D_f,&C_y&=M^2(QV^2-1),\\
 T_N&=\delta^2AB_f,&R_N&=\delta C_y,\\
 J_N&=\delta(Z^2-A-B_f)-C_y,\\
 N_c&=J_N^2-4\bigl(T_N+R_N(A+B_f)\bigr),\\
 N_8&=N_c^2-64T_NR_NZ^2,\\
 P&=N_8+A_1+N_8A_1.
\end{align*}
The literal graphs are \texttt{p11\_positive\_degree277.json} and
\texttt{p11\_positive\_degree257.json}.

\subsection{Three-root separation and the degree-277 lift}

In proof only put $D_y=(QV^2-1)/(Q-1)$, and let
$t=\sqrt{D_\tau}$, $f=\sqrt{D_f}$ and $y=\sqrt{D_y}$ be positive.
For every positive input, independently of any equation or sign of $C$,
\[
 1<D_\tau<D_f<Q,\qquad |V|>1,\qquad D_y>V^2.
\]
The radicand $D_y$ cannot equal either integer $D\in\{D_\tau,D_f\}$:
that would give $Q(D-V^2)=D-1$ with $0<D-1<Q$.
There is also strict root dominance.  Write $b=cT-R$.
If $b=0$, then $V=-c$, $y<c+1$ and $f>t+y$.
If $b\ne0$, then $|V|\ge D_f-c>2f>f+t$, hence $y>f+t$.
These estimates use $I\ge c^2$, $c\ge3$ and
$f>(a+1)I\ge3c^2$.  Thus no signed sum of the three roots is zero.

The universal octic sum norm is
\[
 F_8(x;a,b,c)=
 \bigl[(x^2-a-b-c)^2-4(ab+ac+bc)\bigr]^2-64abcx^2.
\]
It is the product of the eight signed-root factors.  The displayed
$N_8$ is exactly $\delta^4F_8(Z;A,B_f,C_y/\delta)$.
At $N_8=0$, therefore,
\[
 Z/M=\tau_*+f_*+y_*,\qquad
 \tau_*^2=D_\tau,\quad f_*^2=D_f,\quad y_*^2=D_y.
\]
A rational signed sum of three distinct square roots, with no zero
signed sum, forces each root rational.  One exact extraction is as follows:
for signed roots $x_1,x_2,x_3$ with rational nonzero sum $s$ and squares
$a,b,c$, put
\[
 u=(s^2-a-b-c)/2,\quad
 v=(u^2-ab-ac-bc)/(2s),\quad
 x_1=\frac{a(a+u)}{as+v}.
\]
The denominator is $x_1(x_1+x_2)(x_1+x_3)\ne0$; cycle the formula.
Consequently $\tau_*,f_*$ are integers by rational-square integrality,
and $y_*$ is an integer by the preceding Pell descent.

For the $277$-degree construction, $L>f+y$ unconditionally.  Recover
\[
 r=n+L-f_*-y_*>0,\quad f_{\rm old}=|f_*|,\quad
 i_{\rm old}=i,\quad u=|y_*|-V>0.
\]
These are precisely the weighted source equations.  Conversely choose
canonical strong and auxiliary source witnesses first, then $L,W$ and
the transport witness.  The inherited weighted theorem gives
$r>qW-2k>L$, so $n=r-L+f+y>0$.  No stronger divisibility than
$c^2\mid I$ is required.

\subsection{The new centered sign bootstrap}

The centered endpoint cannot assert $r>0$ before proving $V>0$.
The following order resolves that issue.

First recover all three integral roots as above.  Since $M\mid Z$,
reduction modulo $q^3$ and $q-1$ gives
\[
 q^3\mid Y,\qquad q-1\mid B_t.
\]
No transport positivity is used here.  Also $A_1=0$ makes $D_*/U$
integral; classify its \emph{absolute value} as a positive Pell root.
Together with the first Pell root, this gives
\[
 c=\psi_{a+2}(p),\quad k=2\psi_{2XY^2+1}(m),\quad p,m\ge1.
\]
Here $c\ge3$ implies $p\ge2$.  The case $m=1$ would give $k=2$ and
$c=2Y+1<2(a+2)$, a contradiction.  Hence
\[
 c>8XY^3\ge8q^9. \tag{*}
\]

The authentic mask bounds give $|R|<\widehat R$, hence
$|V|<D_fB_0$, while $\delta>D_f$ and $D_f>f$.  Therefore
\[
 0<y-|V|<|V|/\delta<B_0.
\]
In particular $L+V>f+y$ if $V>0$.

Suppose now that $C\le0$.  Set $d=-C\ge0$ and
$B=-B_t=(qK_0+X)d+qF>0$.  Since $K_0\ge1$ and
$(M_C+qM_F)J<2q^2$, direct expansion gives
\[
 qB\ge q^2d+q^2F,\qquad R<qB+q^4+q^2.
\]
Writing $s=Y/q^3>0$, the root equality becomes
\[
 \tau_*+f_*+y_*=
 \ell+k+n+L+V+s+\frac{WB}{q-1}. \tag{**}
\]
If $V>0$, the right side is greater than $t+f+y$, impossible.
If $V<0$, then the integer $cT-R$ is nonpositive, so
$R\ge cT\ge c>8q^9$.  Thus $q^4+q^2<R/2$ and $R<2qB$.
Moreover $|V|<D_fR$, so $L+V-f-y>-2D_fR$, whereas
\[
 \frac{WB}{q-1}>4qD_fB>2D_fR.
\]
Again the right side of $(**)$ exceeds $t+f+y$.
The impossible case $V=0$ was already excluded.  This proves $C>0$.

Only now use the authentic bound $3q<R<q^4-2$.
Together with $(*)$ it gives $V>1$, and therefore
\[
 r=n+L+V-f_*-y_*\ge n+L+V-f-y>n>0.
\]
The old positive strong and auxiliary coordinates are recovered exactly
as before.  The weighted carrier soundness theorem now applies.

For completeness choose canonical strong and auxiliary data before
computing $L,h_*,W$.  Since $h_*>L+k$, the canonical transport variable
satisfies $r>qW-2k>L$.  Set
\[
 n=r-L-V+f+y=r-L+f+u>0.
\]
This gives the displayed centered norm.  The old positive main-carrier
coordinate $\gamma$ is unchanged.  This is an existence equivalence,
not a bijection between every old and new zero set.

Both $N_8$ and $A_1$ are squares modulo four.  Thus
$(1+N_8)(1+A_1)=1$ forces both factors to be $1$, since $-1$ would
require a residual $-2$.  The single equation is therefore equivalent
to the two residual equations.  Under the inherited interfaces, positive
solvability is equivalent to a finite ZFC refutation; negation gives
the stated consistency equivalence.

\subsection{Exact degrees and checks}

For the uncentered construction the exact degrees of
$I,D_f,Q,V,L,W,Z,N_c,N_8,A_1$ are respectively
\[
 5,14,18,18,18,19,24,132,264,13.
\]
The coefficient of
$j^{60}i^{24}\eta^{49}X^{42}Y^{85}\alpha^{16}\gamma$
in the degree-$277$ polynomial is
$2^7 3^8(3\beta)^{60}>0$.

For the centered construction put $z=Z-MV$ in proof only.  The exact
cancellation identity is
\[
 J_N=\delta\{z(z+2MV)-A-B_f\}-M^2(V^2-1).
\]
Here $\deg z=21$, $\deg Z=22$, and the five displayed summand degrees
are $60,61,34,40,44$.  Thus $\deg J_N=61$,
$\deg N_c=122$ and $\deg N_8=244$; the subtracted norm correction
and tail have degrees $102$ and $198$, respectively.  The coefficient of
\[
 j^{52}i^{24}\eta^{49}X^{38}Y^{85}F^4\alpha^4\gamma
\]
is $2^{19}(3\beta)^{52}>0$, proving exact degree $257$ after actual
fixed-port specialization.

Both standard-library builders reproduce the literal graphs byte for
byte.  Both exact checkers pass normally and under \texttt{python -O}:
prefix identities, the eight-factor norm, rational extraction identities,
main-carrier cancellation, centered cancellation, degree heads, selected
coefficients, gate ledgers and fixed coefficient hashes.  Independent
review receipts are kept separately.  These finite checks do not prove
the unbounded native/compiler interfaces or the consistency of ZFC.

\paragraph{Unfinished four-root direction.}
A direct ten-variable extension is not established.  Distinct radicands
and a nonzero full signed sum no longer ensure rationality root by root:
$\sqrt2+\sqrt8-\sqrt{18}+100=100$, with four distinct radicands and
no vanishing full signed sum.  Even a stronger separation argument would
not remove the main-residue versus transport-modulus obstruction from
the earlier attempted merger.  The negative controls in this component
record both issues; no ten-variable endpoint is claimed.
